\section*{Chapitre \ref{ch:rs:point-in-time-data-and-the-biases} --- Données point-in-time et biais}

\begin{solution}{exo:rs:point-in-time-data-and-the-biases:1}
Le cours de clôture du jour $t$ n'est connu qu'à la fin de la séance~; trader à ce même cours utilise donc une information qui n'existe pas encore (sauf si le signal est
calculé à partir de cours de quelques minutes plus tôt et envoyé à la criée de clôture). La première transaction honnête est
l'ouverture de la séance suivante, ou la clôture suivante dans un test clôture à clôture.
\end{solution}

\begin{solution}{exo:rs:point-in-time-data-and-the-biases:2}
$w(r_S - r_F) = 0.15 \times (0.10 + 0.20) = 0.045$~: 4.5 points par an.
\end{solution}

\begin{solution}{exo:rs:point-in-time-data-and-the-biases:3}
En janvier 2009 (quatre mois après), $-403\,000$~; en juin 2009, après l'indice de référence de février, $-321\,000$~; en septembre 2011,
$-434\,000$.
\end{solution}

\begin{solution}{exo:rs:point-in-time-data-and-the-biases:4}
Avec les horodatages~: les cotations horodatées 09:59~:59.5, 10:00~:01.0 et 10:00~:01.8. Avec les heures d'arrivée (09:59~:59.8, 10:00~:01.3, 10:00~:02.1)~: la
première cotation pour les deux premières décisions et la seconde pour la troisième. La jointure sur les horodatages permet à la seconde décision d'utiliser une cotation
0.3 seconde avant qu'elle n'existe chez la société.
\end{solution}

\begin{solution}{exo:rs:point-in-time-data-and-the-biases:5}
La surprise de la première publication est la surprise sur les données finales moins la révision~; sa variance vaut donc $90^2 + 106^2 = 19\,336$ (en
milliers d'emplois au carré), dont $11\,236$, soit 58~\%, est du bruit de révision. Un modèle calé sur les données finales voit un signal bien plus net que celui
sur lequel le marché a réellement opéré.
\end{solution}

\begin{solution}{exo:rs:point-in-time-data-and-the-biases:6}
L'historique antérieur à l'entrée d'un fonds n'a jamais été disponible en temps réel et il est sélectionné~: seuls les fonds ayant de bons rendements passés rejoignent la base, si bien que
leurs années rétroremplies relèvent la moyenne de la base (\omterm{def:rs:point-in-time-data-and-the-biases:vintage}{biais de remplissage rétroactif}). Pour le supprimer, on n'utilise les rendements de chaque fonds qu'à partir de sa
date d'entrée dans la base (son \omterm{def:rs:point-in-time-data-and-the-biases:bitemporal}{temps de connaissance}), ou l'on écarte une période initiale fixe de chaque fonds.
\end{solution}

\begin{solution}{exo:rs:point-in-time-data-and-the-biases:7}
Seuil $-1.2$~: 3.5~\% des titres radiés par an, biais de 7.0 points~; seuil $-2.0$~: 1.4~\% et 3.8 points. Rendement de radiation nul~: 4.5
points~; $-100~\%$~: 6.7 points. Le seuil déplace davantage le biais, car il fixe le nombre de titres qui défaillent et la durée pendant laquelle
ils déclinent~; le rendement de radiation n'affecte qu'un mois de chaque défaillance.
\end{solution}

\begin{solution}{exo:rs:point-in-time-data-and-the-biases:8}
Trois biais à la fois. \omterm{def:rs:point-in-time-data-and-the-biases:survivorship}{Biais de survie}~: les membres actuels excluent les sociétés qui ont fait défaut ou se sont contractées au point de sortir de l'indice depuis 2000.
\omterm{def:rs:point-in-time-data-and-the-biases:lookahead}{Biais d'anticipation} sur l'appartenance à l'indice~: beaucoup de membres actuels y ont été ajoutés parce qu'ils avaient progressé, un fait inconnu en 2000. \omterm{def:rs:point-in-time-data-and-the-biases:vintage}{Retraitements}
et corrections des fournisseurs~: les fondamentaux d'aujourd'hui ne sont pas ceux publiés à l'époque. Le test exige l'appartenance à l'indice à chaque
date, les sociétés radiées avec leurs rendements de radiation et des fondamentaux point-in-time.
\end{solution}

\begin{solution}{pb:rs:point-in-time-data-and-the-biases:1}
\textbf{1.} $1 - e^{-1.6} = 79.8~\%$. \textbf{2.} 2.2~\% des titres par an. \textbf{3.} 7.7~\% par an. \textbf{4.} 12.9~\%.
\textbf{5.} 5.2 points. \textbf{6.} 20.9~\%. \textbf{7.} $r_F = -11.1~\%$, $r_S = 12.6~\%$ par an. \textbf{8.} $0.209 \times (12.6 +
11.1) = 5.0$ points~; la question~5 calcule la moyenne mois par mois sur un nombre de survivants variable (peu au début, beaucoup à la fin) au lieu
de regrouper les couples titre-mois. \textbf{9.} 4.5 points. \textbf{10.} 6.7 points. \textbf{11.} 3.5~\% par an, 7.0 points. \textbf{12.} 1.4~\%
par an, 3.8 points. \textbf{13.} Davantage de titres atteignent le seuil, donc davantage de couples titre-mois relèvent des défaillances ($w$ augmente), tandis que
le recul de chaque défaillance reste prononcé. \textbf{14.} Un profit élevé~: les titres tombés qui ont fini par faire faillite manquent, et ceux qui
restent sont par construction ceux qui se sont redressés~; le véritable résultat peut être de signe contraire. \textbf{15.} À partir d'une base qui
conserve les titres radiés et leurs rendements de radiation, avec des univers arrêtés à chaque date. \textbf{16.} Une estimation par motif
de radiation, comme le $-30~\%$ que Shumway a récupéré pour les radiations liées à la performance~; un rendement nul est optimiste, car les radiations liées à la performance sont
des surprises que le dernier cours coté ne reflète pas. \textbf{17.} Les titres en défaillance sont petits au moment où ils font défaut, donc un
univers pondéré par la capitalisation leur accorde peu de poids et souffre moins~; la pondération égale, courante en recherche, souffre le plus.
\textbf{18.} \emph{Résultat chiffré~:} l'univers des seuls survivants surestime le rendement pondéré également de 5.2 points par an lorsque 2.2~\%
des titres sont radiés chaque année, et le mois de radiation en explique 0.7 (le biais reste de 4.5 points avec un rendement de radiation
nul). \textbf{19.} \omterm{def:rs:point-in-time-data-and-the-biases:vintage}{Biais de remplissage rétroactif} dans les bases de fonds et de fournisseurs, et test sur les membres actuels de l'indice. \textbf{20.} Elle recense les
sociétés qui survivront, ce que personne ne connaissait au début de l'échantillon.
\end{solution}

\begin{solution}{iq:rs:point-in-time-data-and-the-biases:1}
Étudier uniquement les titres qui existent à la fin de l'échantillon. Dans un backtest de sélection d'actions construit sur les sociétés cotées aujourd'hui,
les défaillances manquent, de sorte que les stratégies qui achètent des actions bon marché, en difficulté ou petites paraissent bien meilleures qu'elles ne l'étaient, et les rendements moyens
sont gonflés de plusieurs points par an dans les univers à pondération égale.
\iqlookfor{que les titres manquants sont les perdants, et la correction~: un univers arrêté à chaque date avec les radiations.}
\end{solution}

\begin{solution}{iq:rs:point-in-time-data-and-the-biases:2}
L'estimation disponible à chaque date de décision~: le millésime publié à cette date (première ou ultérieure publication), et non le chiffre révisé
d'aujourd'hui~; le modèle doit aussi être calé sur ces valeurs en temps réel, puisqu'il les recevra en direct.
\iqlookfor{millésimes en temps réel~; la distinction entre \omterm{def:rs:point-in-time-data-and-the-biases:bitemporal}{temps de validité} et \omterm{def:rs:point-in-time-data-and-the-biases:bitemporal}{temps de connaissance}.}
\end{solution}

\begin{solution}{iq:rs:point-in-time-data-and-the-biases:3}
Clé (société, champ, exercice comptable, \omterm{def:rs:point-in-time-data-and-the-biases:bitemporal}{temps de connaissance}) associée à la valeur, où le \omterm{def:rs:point-in-time-data-and-the-biases:bitemporal}{temps de connaissance} est l'heure du dépôt ou de la livraison par le fournisseur~;
ne jamais écraser. Une requête prend (société, champ, période, instant as-of) et renvoie la dernière ligne connue à cet instant~; une requête
d'instantané renvoie chaque période telle qu'elle était connue alors.
\iqlookfor{deux colonnes temporelles, stockage en ajout seul, sémantique as-of.}
\end{solution}

\begin{solution}{iq:rs:point-in-time-data-and-the-biases:4}
Perturbation~: pour chaque instant de décision, remplacer par du bruit chaque valeur d'entrée dont le \omterm{def:rs:point-in-time-data-and-the-biases:bitemporal}{temps de connaissance} est postérieur, et vérifier que les
caractéristiques à cette décision sont inchangées. Ajouter aussi un garde-fou dans la couche d'accès aux données, qui lève une erreur pour toute date de connaissance postérieure, et
une comparaison du backtest exécuté sur les millésimes conservés avec le même calcul sur les données actuelles (un écart important signale une sensibilité aux
révisions).
\iqlookfor{un test automatisé, et non une revue de code.}
\end{solution}

\begin{solution}{iq:rs:point-in-time-data-and-the-biases:5}
\omterm{def:rs:point-in-time-data-and-the-biases:survivorship}{Biais de survie} (fonds disparus écartés), remplissage rétroactif (fonds qui entrent dans la base avec leur bon passé), autosélection (fonds qui publient quand ils vont bien
et cessent de publier quand ils vont mal), et valorisations périmées ou lissées pour les positions peu liquides, qui abaissent la volatilité mesurée.
\iqlookfor{plusieurs biais distincts et leur sens.}
\end{solution}

\begin{solution}{iq:rs:point-in-time-data-and-the-biases:6}
Agrégée sur les couples titre-mois, la moyenne de l'univers vaut $(1 - w)r_S + wr_F$~; la moyenne des survivants vaut $r_S$~; l'écart vaut $w(r_S
- r_F)$. En pratique, l'écart $r_S - r_F$ est dominé par les mois pendant lesquels les titres en défaillance déclinent, et non par le mois de radiation~: dans la
simulation du chapitre, 4.5 des 5.2 points subsistent avec un rendement de radiation nul.
\iqlookfor{l'identité, et d'où vient l'écart.}
\end{solution}
